\documentclass[10pt,a4paper]{amsart}
\usepackage{amsmath,amssymb,amsthm,graphicx}
\usepackage{verbatim}
\newenvironment{reference}{\comment}{\endcomment}
%\newenvironment{reference}{}{}
\newenvironment{slogan}{\comment}{\endcomment}
\newenvironment{history}{\comment}{\endcomment}

% For commutative diagrams we use Xy-pic
\usepackage[all]{xy}

% We use 2cell for 2-commutative diagrams.
\xyoption{2cell}
% Theorem environments.
%
\theoremstyle{plain}
\newtheorem{theorem}[subsection]{Theorem}
\newtheorem{proposition}[subsection]{Proposition}
\newtheorem{lemma}[subsection]{Lemma}

\theoremstyle{definition}
\newtheorem{definition}[subsection]{Definition}
\newtheorem{example}[subsection]{Example}
\newtheorem{exercise}[subsection]{Exercise}
\newtheorem{situation}[subsection]{Situation}

\theoremstyle{remark}
\newtheorem{remark}[subsection]{Remark}
\newtheorem{remarks}[subsection]{Remarks}

\numberwithin{equation}{subsection}

% Macros
%
\def\lim{\mathop{\mathrm{lim}}\nolimits}
\def\colim{\mathop{\mathrm{colim}}\nolimits}
\def\Spec{\mathop{\mathrm{Spec}}}
\def\Hom{\mathop{\mathrm{Hom}}\nolimits}
\def\Ext{\mathop{\mathrm{Ext}}\nolimits}
\def\SheafHom{\mathop{\mathcal{H}\!\mathit{om}}\nolimits}
\def\SheafExt{\mathop{\mathcal{E}\!\mathit{xt}}\nolimits}
\def\Sch{\mathit{Sch}}
\def\Mor{\mathop{\mathrm{Mor}}\nolimits}
\def\Ob{\mathop{\mathrm{Ob}}\nolimits}
\def\Sh{\mathop{\mathit{Sh}}\nolimits}
\def\NL{\mathop{N\!L}\nolimits}
\def\CH{\mathop{\mathrm{CH}}\nolimits}
\def\proetale{{pro\text{-}\acute{e}tale}}
\def\etale{{\acute{e}tale}}
\def\QCoh{\mathit{QCoh}}
\def\Ker{\mathop{\mathrm{Ker}}}
\def\Im{\mathop{\mathrm{Im}}}
\def\Coker{\mathop{\mathrm{Coker}}}
\def\Coim{\mathop{\mathrm{Coim}}}

% Boxtimes
%
\DeclareMathSymbol{\boxtimes}{\mathbin}{AMSa}{"02}

%
% Macros for moduli stacks/spaces
%
\def\QCohstack{\mathcal{QC}\!\mathit{oh}}
\def\Cohstack{\mathcal{C}\!\mathit{oh}}
\def\Spacesstack{\mathcal{S}\!\mathit{paces}}
\def\Quotfunctor{\mathrm{Quot}}
\def\Hilbfunctor{\mathrm{Hilb}}
\def\Curvesstack{\mathcal{C}\!\mathit{urves}}
\def\Polarizedstack{\mathcal{P}\!\mathit{olarized}}
\def\Complexesstack{\mathcal{C}\!\mathit{omplexes}}
% \Pic is the operator that assigns to X its picard group, usage \Pic(X)
% \Picardstack_{X/B} denotes the Picard stack of X over B
% \Picardfunctor_{X/B} denotes the Picard functor of X over B
\def\Pic{\mathop{\mathrm{Pic}}\nolimits}
\def\Picardstack{\mathcal{P}\!\mathit{ic}}
\def\Picardfunctor{\mathrm{Pic}}
\def\Deformationcategory{\mathcal{D}\!\mathit{ef}}
\usepackage{cleveref}
\crefname{originalreference}{original reference}{original references}
\Crefname{originalreference}{Original reference}{Original references}
\makeatletter
\newlabel{lemma-base-change-tensor}{{0A1D}{source}{}{sourcepointer-1}{}}
\newlabel{lemma-base-change-tensor@cref}{{[originalreference][0][]0A1D}{[1][0][]source}{}{}{}}
\makeatother
\crefname{originalreference}{original reference}{original references}
\Crefname{originalreference}{Original reference}{Original references}
\makeatletter
\newlabel{lemma-perfect-on-limit}{{09RJ}{source}{}{sourcepointer-2}{}}
\newlabel{lemma-perfect-on-limit@cref}{{[originalreference][0][]09RJ}{[1][0][]source}{}{}{}}
\makeatother
\crefname{originalreference}{original reference}{original references}
\Crefname{originalreference}{Original reference}{Original references}
\makeatletter
\newlabel{lemma-tensor-perfect}{{09J5}{source}{}{sourcepointer-3}{}}
\newlabel{lemma-tensor-perfect@cref}{{[originalreference][0][]09J5}{[1][0][]source}{}{}{}}
\makeatother
\crefname{originallemma}{original lemma}{original lemmas}
\Crefname{originallemma}{Original lemma}{Original lemmas}
\makeatletter
\newlabel{spaces-limits-lemma-descend-finite-presentation}{{07SK}{source}{}{sourcepointer-4}{}}
\newlabel{spaces-limits-lemma-descend-finite-presentation@cref}{{[originallemma][0][]07SK}{[1][0][]source}{}{}{}}
\makeatother
\crefname{originallemma}{original lemma}{original lemmas}
\Crefname{originallemma}{Original lemma}{Original lemmas}
\makeatletter
\newlabel{spaces-limits-lemma-descend-flat}{{08K0}{source}{}{sourcepointer-5}{}}
\newlabel{spaces-limits-lemma-descend-flat@cref}{{[originallemma][0][]08K0}{[1][0][]source}{}{}{}}
\makeatother
\crefname{originallemma}{original lemma}{original lemmas}
\Crefname{originallemma}{Original lemma}{Original lemmas}
\makeatletter
\newlabel{spaces-limits-lemma-descend-modules-finite-presentation}{{07V7}{source}{}{sourcepointer-6}{}}
\newlabel{spaces-limits-lemma-descend-modules-finite-presentation@cref}{{[originallemma][0][]07V7}{[1][0][]source}{}{}{}}
\makeatother
\crefname{originallemma}{original lemma}{original lemmas}
\Crefname{originallemma}{Original lemma}{Original lemmas}
\makeatletter
\newlabel{spaces-limits-lemma-eventually-proper-support}{{08K2}{source}{}{sourcepointer-7}{}}
\newlabel{spaces-limits-lemma-eventually-proper-support@cref}{{[originallemma][0][]08K2}{[1][0][]source}{}{}{}}
\makeatother
\title{Tensoring a perfect complex with a nice complex pushes forward to a perfect complex with base-change compatible formation}
\author{The Stacks Project Authors}
\date{stacks-project:spaces-perfect.tex:local-snapshot-20261012; GNU FDL 1.2 or later}
\begin{document}
\maketitle
\noindent\textit{Editorial note.} All delivered mathematical text is the author's own, copied from the pinned Stacks Project chapter file spaces-perfect.tex; no statement or proof step was rewritten, reordered or removed.\par\medskip
\noindent\textit{Editorial note.} The counted proof is the author's complete argument: the proof is complete: the Noetherian case is quoted and the general case descends along a filtered colimit\par\medskip
\noindent\textit{Editorial note.} The label lemma-base-change-tensor-perfect names the counted statement, and the proof environment immediately following it in the source is delivered verbatim as its proof.\par\medskip
\noindent\textit{Editorial note.} The counted statement is carried by the author's own lemma environment in the chapter source.\par\medskip
\noindent\textit{Allowed standard prerequisites (editorial).}\begin{itemize}
\item the base-change compatibility of the total direct image of a tensor product (base-change-tensor lemma)
\item the tensor-perfect lemma over Noetherian bases
\item Noetherian approximation: descent of finite presentation, modules, flatness and properness of support along filtered colimits of rings
\item the perfect-on-limit lemma descending perfect objects
\end{itemize}\medskip
\ifdefined\hypertarget\hypertarget{sourcepointer-1}{}\fi
\noindent\textit{Source reference.} Original reference 0A1D: sources/stacks-spaces-perfect/pinned-original-source.tex, source lines 5625--5625. Prerequisite pointer invoked at proof line(s) 6321; it names a standard fact of the same chapter that the counted argument uses, and every step of the counted proof built on it is retained verbatim in the delivered body.\par\medskip
\ifdefined\hypertarget\hypertarget{sourcepointer-2}{}\fi
\noindent\textit{Source reference.} Original reference 09RJ: sources/stacks-spaces-perfect/pinned-original-source.tex, source lines 6229--6229. Prerequisite pointer invoked at proof line(s) 6347; it names a standard fact of the same chapter that the counted argument uses, and every step of the counted proof built on it is retained verbatim in the delivered body.\par\medskip
\ifdefined\hypertarget\hypertarget{sourcepointer-3}{}\fi
\noindent\textit{Source reference.} Original reference 09J5: sources/stacks-spaces-perfect/pinned-original-source.tex, source lines 5794--5794. Prerequisite pointer invoked at proof line(s) 6324; it names a standard fact of the same chapter that the counted argument uses, and every step of the counted proof built on it is retained verbatim in the delivered body.\par\medskip
\ifdefined\hypertarget\hypertarget{sourcepointer-4}{}\fi
\noindent\textit{Source reference.} Original lemma 07SK: sources/stacks-chow/tags.txt, source lines 9905--9905. Prerequisite pointer invoked at proof line(s) 6332; it names a standard fact of a sibling chapter that the counted argument uses, and every step of the counted proof built on it is retained verbatim in the delivered body.\par\medskip
\ifdefined\hypertarget\hypertarget{sourcepointer-5}{}\fi
\noindent\textit{Source reference.} Original lemma 08K0: sources/stacks-chow/tags.txt, source lines 10868--10868. Prerequisite pointer invoked at proof line(s) 6343; it names a standard fact of a sibling chapter that the counted argument uses, and every step of the counted proof built on it is retained verbatim in the delivered body.\par\medskip
\ifdefined\hypertarget\hypertarget{sourcepointer-6}{}\fi
\noindent\textit{Source reference.} Original lemma 07V7: sources/stacks-chow/tags.txt, source lines 9997--9997. Prerequisite pointer invoked at proof line(s) 6336; it names a standard fact of a sibling chapter that the counted argument uses, and every step of the counted proof built on it is retained verbatim in the delivered body.\par\medskip
\ifdefined\hypertarget\hypertarget{sourcepointer-7}{}\fi
\noindent\textit{Source reference.} Original lemma 08K2: sources/stacks-chow/tags.txt, source lines 10870--10870. Prerequisite pointer invoked at proof line(s) 6346; it names a standard fact of a sibling chapter that the counted argument uses, and every step of the counted proof built on it is retained verbatim in the delivered body.\par\medskip
\begin{lemma}
\label{lemma-base-change-tensor-perfect}
Let $S$ be a scheme. Let $f : X \to Y$ be a morphism of finite presentation
between algebraic spaces over $S$. Let $E \in D(\mathcal{O}_X)$ be a perfect
object. Let $\mathcal{G}^\bullet$ be a bounded complex of finitely presented
$\mathcal{O}_X$-modules, flat over $Y$, with support proper over $Y$. Then
$$
K = Rf_*(E \otimes_{\mathcal{O}_X}^\mathbf{L} \mathcal{G}^\bullet)
$$
is a perfect object of $D(\mathcal{O}_Y)$ and its formation
commutes with arbitrary base change.
\end{lemma}

\begin{proof}
The statement on base change is Lemma \ref{lemma-base-change-tensor}.
Thus it suffices to show that $K$ is a perfect object. If $Y$ is
Noetherian, then this follows from
Lemma \ref{lemma-tensor-perfect}.
We will reduce to this case by Noetherian approximation.
We encourage the reader to skip the rest of this proof.

\medskip\noindent
The question is local on $Y$, hence we may assume $Y$ is affine.
Say $Y = \Spec(R)$. We write $R = \colim R_i$ as a filtered colimit
of Noetherian rings $R_i$. By Limits of Spaces, Lemma
\ref{spaces-limits-lemma-descend-finite-presentation}
there exists an $i$ and an algebraic space $X_i$ of finite presentation
over $R_i$ whose base change to $R$ is $X$. By
Limits of Spaces, Lemma
\ref{spaces-limits-lemma-descend-modules-finite-presentation}
we may assume after increasing $i$, that there exists a
bounded complex of finitely
presented $\mathcal{O}_{X_i}$-modules $\mathcal{G}_i^\bullet$ whose
pullback to $X$ is $\mathcal{G}^\bullet$. After increasing $i$
we may assume $\mathcal{G}_i^n$ is flat over $R_i$, see
Limits of Spaces, Lemma
\ref{spaces-limits-lemma-descend-flat}.
After increasing $i$ we may assume the support of $\mathcal{G}_i^n$
is proper over $R_i$, see
Limits of Spaces, Lemma \ref{spaces-limits-lemma-eventually-proper-support}.
Finally, by Lemma \ref{lemma-perfect-on-limit}
we may, after increasing $i$, assume there exists a perfect
object $E_i$ of $D(\mathcal{O}_{X_i})$ whose pullback to
$X$ is $E$. By Lemma \ref{lemma-tensor-perfect} we have that
$K_i =
Rf_{i, *}(E_i \otimes_{\mathcal{O}_{X_i}}^\mathbf{L} \mathcal{G}_i^\bullet)$
is perfect on $\Spec(R_i)$ where $f_i : X_i \to \Spec(R_i)$ is
the structure morphism. By the base change result
(Lemma \ref{lemma-base-change-tensor}) the pullback of $K_i$
to $Y = \Spec(R)$ is $K$ and we conclude.
\end{proof}

\end{document}
